%==============================================================================
% CAPÍTULO 6 — REDES NEURAIS ARTIFICIAIS
% Parte II: Teoria e Modelos de Deep Learning
%==============================================================================

\chapter{Redes Neurais Artificiais}
\label{cap:redes-neurais}

\niveltres

\begin{quote}
\itshape\small
``Neural networks are a new technique for development of decision support systems in dentistry.'' --- \citeonline{ODT-002}, Journal of Dentistry, 1998.
\end{quote}

Este capítulo estabelece os fundamentos matemáticos e computacionais das redes neurais
artificiais (RNAs), a espinha dorsal do deep learning moderno. Partimos do modelo mais
simples --- o perceptron --- e progredimos até as redes multicamadas com retropropagação,
cobrindo funções de ativação, funções de perda, otimizadores e técnicas de regularização.
Ao final, implementaremos uma rede MLP (Multi-Layer Perceptron) utilizando o módulo
\texttt{odontoia.models.cnn} para classificação de condições dentárias.

\notaexplicativa{Este capítulo é intensivo em notação matemática. Recomendamos que o
leitor tenha familiaridade com álgebra linear básica (operações matriciais, derivação
de funções compostas). O Apêndice A contém um glossário com todos os termos técnicos.}

\section{Objetivos de Aprendizagem}

Ao final deste capítulo, você será capaz de:

\subsection{Entender o Funcionamento do Perceptron}
Você aprenderá como o neurônio artificial de Rosenblatt processa entradas ponderadas e produz saídas, formando a base de todas as redes neurais modernas.

\subsection{Dominar as Funções de Ativação}
Exploraremos sigmoide, ReLU, tanh e suas variantes, compreendendo como cada uma influencia o aprendizado da rede.

\subsection{Implementar uma MLP para Classificação Odontológica}
Ao final, você implementará uma rede multicamadas capaz de classificar condições dentárias em radiografias.

\subsection{Compreender o Algoritmo de Retropropagação}
Você entenderá como o gradiente da perda é calculado camada por camada, capacitando-o a diagnosticar e corrigir problemas de treinamento.

%-----------------------------------------------------------------------------
\section{O Neurônio Artificial}
\label{sec:neuronio-artificial}

\subsection{Inspiração Biológica}

O neurônio biológico recebe sinais elétricos através de dendritos, integra-os no corpo
celular (soma) e, caso o potencial de membrana atinja um limiar, dispara um potencial
de ação que se propaga pelo axônio até as sinapses. Este mecanismo binário de
``tudo-ou-nada'' inspirou McCulloch e Pitts (1943) a propor o primeiro modelo
matemático de neurônio artificial.

No modelo computacional, o neurônio artificial é uma unidade de processamento que:
\begin{enumerate}
  \item Recebe $n$ entradas (features) $\mathbf{x} = [x_1, x_2, \ldots, x_n]^T$
  \item Pondera cada entrada por um peso sináptico $w_i$
  \item Soma os sinais ponderados e adiciona um bias $b$
  \item Aplica uma função de ativação não-linear $f(\cdot)$
  \item Produz a saída $\hat{y} = f\left(\sum_{i=1}^{n} w_i x_i + b\right)$
\end{enumerate}

\subsection{O Perceptron de Rosenblatt}

\subsubsection{O Algoritmo do Perceptron}

Frank Rosenblatt (1958) implementou o perceptron --- o primeiro neurônio artificial
treinável --- no hardware Mark I Perceptron. O algoritmo de aprendizado do perceptron
é notavelmente simples:

\begin{equation}
  \hat{y} = \text{sign}\left(\sum_{i=1}^{n} w_i x_i + b\right) = \begin{cases}
    +1, & \text{se } \mathbf{w}^T\mathbf{x} + b \geq 0 \\
    -1, & \text{caso contrário}
  \end{cases}
  \label{eq:perceptron}
\end{equation}

A regra de atualização dos pesos para cada exemplo de treinamento $(\mathbf{x}^{(j)}, y^{(j)})$:

\begin{equation}
  w_i \leftarrow w_i + \eta \cdot (y^{(j)} - \hat{y}^{(j)}) \cdot x_i^{(j)}
  \label{eq:perceptron-update}
\end{equation}

Onde $\eta$ é a taxa de aprendizado (\textit{learning rate}) e $(y^{(j)} - \hat{y}^{(j)})$
é o erro de classificação. O algoritmo converge em um número finito de passos apenas
quando os dados são linearmente separáveis --- o que motivou a criação de arquiteturas
mais complexas.

\subsubsection{Limitações do Perceptron}

\begin{table}[H]
  \centering
  \caption{Limitações do perceptron simples e suas soluções históricas}
  \label{tab:limitacoes-perceptron}
  \begin{tabularx}{\textwidth}{p{3.5cm}XX}
    \toprule
    \textbf{Limitação} & \textbf{Problema Clássico} & \textbf{Solução} \\
    \midrule
    Separação linear & XOR (Minsky \& Papert, 1969) & Camadas ocultas + não-linearidades \\
    Apenas 2 classes & Classificação multiclasse odontológica & Softmax + one-vs-all \\
    Gradiente binário & Função sign não-diferenciável & Funções de ativação contínuas \\
    Sem probabilidades & Diagnóstico incerto & Sigmoid/softmax como saída \\
    \bottomrule
  \end{tabularx}
\end{table}

\citacaoativa{doi-1016-j-jds-2020-06-019}{10.1016/j.jds.2020.06.019}{A transição de sistemas especialistas baseados
em regras para redes neurais e deep learning marca a principal mudança de paradigma na
odontologia computacional das últimas duas décadas (Khanagar et al., 2021).}

%-----------------------------------------------------------------------------
\section{Funções de Ativação}
\label{sec:funcoes-ativacao}

A função de ativação introduz a \textbf{não-linearidade} essencial que permite às RNAs
aproximarem funções arbitrariamente complexas. Sem ela, uma rede de $L$ camadas seria
equivalente a uma única transformação linear, independentemente da profundidade.

\subsection{Sigmoide ($\sigma$)}

A função sigmoide ``comprime'' qualquer valor real para o intervalo $(0, 1)$,
tornando-a útil para saídas probabilísticas binárias.

\begin{equation}
  \sigma(z) = \frac{1}{1 + e^{-z}}, \quad \sigma'(z) = \sigma(z)(1 - \sigma(z))
  \label{eq:sigmoide}
\end{equation}

\textbf{Vantagens:} interpretação probabilística natural; derivada simples.
\textbf{Desvantagens:} saturação para $|z| > 5$ causa gradientes próximos de zero
(\textit{vanishing gradient}); saída não-centrada em zero (média $\approx 0.5$);
custo computacional da exponencial.

Na odontologia, a sigmoide é comum na camada de saída de classificadores binários
de lesões (maligno vs. benigno), como no trabalho de \citeonline{DL-ODT-020}.

\subsection{Tangente Hiperbólica ($\tanh$)}

\begin{equation}
  \tanh(z) = \frac{e^z - e^{-z}}{e^z + e^{-z}}, \quad \tanh'(z) = 1 - \tanh^2(z)
  \label{eq:tanh}
\end{equation}

A $\tanh$ é uma versão escalonada e centrada em zero da sigmoide (saída entre $-1$
e $1$), mitigando parcialmente o problema de convergência lenta. Contudo, ainda sofre
de saturação nos extremos.

\subsection{ReLU (\textit{Rectified Linear Unit})}

\subsubsection{Definição e Propriedades}

A ReLU é a função de ativação mais utilizada em arquiteturas CNN modernas, inclusive
nas aplicações odontológicas revisadas por \citeonline{DL-ODT-004}.

\begin{equation}
  \text{ReLU}(z) = \max(0, z), \quad \text{ReLU}'(z) = \begin{cases}
    1, & z > 0 \\
    0, & z \leq 0
  \end{cases}
  \label{eq:relu}
\end{equation}

\textbf{Vantagens:} computacionalmente barata (apenas uma comparação e um threshold);
não satura para valores positivos (gradiente constante = $1$); acelera convergência
do SGD em $6\times$ comparado à sigmoide/$\tanh$ \cite{lecun2015deep}.

\textbf{Desvantagens:} ``neurônios mortos'' --- quando $z \leq 0$, o gradiente é zero
e o neurônio para de aprender permanentemente. Este problema pode afetar até 40\%
das unidades em arquiteturas profundas.

\subsubsection{Variantes da ReLU}

\textbf{Variantes que mitigam neurônios mortos:}
\begin{itemize}
  \item \textbf{Leaky ReLU:} $f(z) = \max(0.01z, z)$ --- pequeno gradiente para $z < 0$
  \item \textbf{Parametric ReLU (PReLU):} $f(z) = \max(\alpha z, z)$, com $\alpha$
        aprendido durante o treinamento
  \item \textbf{ELU (\textit{Exponential Linear Unit}):}
        $f(z) = \begin{cases} z, & z > 0 \\ \alpha(e^z - 1), & z \leq 0 \end{cases}$
  \item \textbf{GELU (\textit{Gaussian Error Linear Unit}):} utilizada nos
        Transformers (ViT, BERT), aproxima a ReLU com suavidade gaussiana
\end{itemize}

\subsection{Softmax}

A softmax é usada exclusivamente na \textbf{camada de saída} para classificação
multiclasse, convertendo logits em uma distribuição de probabilidade.

\begin{equation}
  \text{softmax}(z_i) = \frac{e^{z_i}}{\sum_{j=1}^{K} e^{z_j}}, \quad
  \sum_{i=1}^{K} \text{softmax}(z_i) = 1
  \label{eq:softmax}
\end{equation}

No contexto odontológico, uma classificação de 5 tipos de lesões bucais (carcinoma,
leucoplasia, eritroplasia, líquen plano, normal) utiliza softmax com $K=5$
(\citeonline{ODT-013}). A classe predita é $\arg\max_i \text{softmax}(z_i)$.

\begin{table}[H]
  \centering
  \caption{Comparação de funções de ativação para tarefas odontológicas típicas}
  \label{tab:ativacao-odonto}
  \begin{tabularx}{\textwidth}{lcccX}
    \toprule
    \textbf{Função} & \textbf{Saída} & \textbf{Camada} & \textbf{Uso Típico} & \textbf{Exemplo Odontológico} \\
    \midrule
    ReLU & $[0, \infty)$ & Oculta/CNN & Padrão em CNNs & Detecção de cáries (Cantu 2023) \\
    $\tanh$ & $[-1, 1]$ & Oculta (RNN/LSTM) & Dados normalizados & Predição de sobrevida OSCC \\
    Sigmoide & $(0, 1)$ & Saída binária & 2 classes & Maligno vs. benigno \\
    Softmax & $(0,1)^K$ & Saída multiclasse & $K \geq 3$ classes & 5 classes de lesões bucais \\
    \bottomrule
  \end{tabularx}
\end{table}

%-----------------------------------------------------------------------------
\section{Funções de Perda}
\label{sec:funcoes-perda}

A função de perda $\mathcal{L}(\mathbf{\hat{y}}, \mathbf{y})$ quantifica a discrepância
entre as predições do modelo e os valores reais. O treinamento consiste em minimizar
esta função ajustando os parâmetros $\mathbf{W}, \mathbf{b}$.

\subsection{Erro Quadrático Médio (MSE)}

Para problemas de \textbf{regressão} (predição de valores contínuos):

\begin{equation}
  \mathcal{L}_{\text{MSE}} = \frac{1}{m} \sum_{i=1}^{m} (\hat{y}^{(i)} - y^{(i)})^2
  \label{eq:mse}
\end{equation}

Na odontologia, a MSE pode ser usada na predição da profundidade de sondagem
periodontal (valores em mm) ou no treinamento de autoencoders para reconstrução
de imagens CBCT com artefatos reduzidos (\citeonline{DL-ODT-015}).

\subsection{Entropia Cruzada Binária (BCE)}

Para \textbf{classificação binária}:

\begin{equation}
  \mathcal{L}_{\text{BCE}} = -\frac{1}{m} \sum_{i=1}^{m} \left[
    y^{(i)} \log(\hat{y}^{(i)}) + (1 - y^{(i)}) \log(1 - \hat{y}^{(i)})
  \right]
  \label{eq:bce}
\end{equation}

\notaexplicativa{A entropia cruzada é derivada do princípio de máxima verossimilhança.
Quando combinada com a ativação sigmoide, o gradiente simplifica-se para
$\hat{y} - y$, evitando o problema de \textit{vanishing gradient} da MSE com sigmoide.}

\subsection{Entropia Cruzada Categórica}

Para \textbf{classificação multiclasse} com $K$ categorias (ex.: 32 dentes na
numeração FDI):

\begin{equation}
  \mathcal{L}_{\text{CE}} = -\frac{1}{m} \sum_{i=1}^{m} \sum_{k=1}^{K}
    y_k^{(i)} \log(\hat{y}_k^{(i)})
  \label{eq:ce}
\end{equation}

\subsection{Dice Loss (Segmentação)}

Específica para tarefas de segmentação (Capítulo 8), onde há forte desbalanceamento
entre pixels de primeiro plano (lesão/dente) e fundo:

\begin{equation}
  \mathcal{L}_{\text{Dice}} = 1 - \frac{2 \sum_i p_i g_i}{\sum_i p_i^2 + \sum_i g_i^2}
  \label{eq:dice-loss}
\end{equation}

Onde $p_i$ é a probabilidade predita e $g_i$ é o ground truth (0 ou 1) para o pixel $i$.

%-----------------------------------------------------------------------------
\section{Backpropagation}
\label{sec:backpropagation}

O algoritmo de retropropagação (\textit{backpropagation}), popularizado por
Rumelhart, Hinton e Williams (1986), é o mecanismo que permite treinar redes
multicamadas calculando gradientes de forma eficiente via regra da cadeia.

\subsection{Intuição em Quatro Passos}

\begin{enumerate}
  \item \textbf{Forward Pass:} A entrada $\mathbf{x}$ propaga-se camada por camada,
        produzindo uma predição $\mathbf{\hat{y}}$.
  \item \textbf{Cálculo da Perda:} $\mathcal{L}(\mathbf{\hat{y}}, \mathbf{y})$
        quantifica o erro.
  \item \textbf{Backward Pass:} O gradiente da perda em relação a cada parâmetro é
        calculado da última camada para a primeira, usando a regra da cadeia:
        \begin{equation}
          \frac{\partial \mathcal{L}}{\partial w_{ij}^{(l)}} =
          \frac{\partial \mathcal{L}}{\partial a_j^{(l)}} \cdot
          \frac{\partial a_j^{(l)}}{\partial z_j^{(l)}} \cdot
          \frac{\partial z_j^{(l)}}{\partial w_{ij}^{(l)}}
          \label{eq:chain-rule}
        \end{equation}
        Onde $a_j^{(l)}$ é a ativação e $z_j^{(l)} = \sum_i w_{ij}^{(l)} a_i^{(l-1)} + b_j^{(l)}$.
  \item \textbf{Atualização dos Parâmetros:} Gradiente descendente:
        $w_{ij}^{(l)} \leftarrow w_{ij}^{(l)} - \eta \frac{\partial \mathcal{L}}{\partial w_{ij}^{(l)}}$
\end{enumerate}

\subsection{Exemplo Numérico: XOR com MLP (2-2-1)}

Para fixar o conceito, apresentamos o cálculo completo para uma rede com uma camada
oculta resolvendo o problema XOR:

\subsubsection*{Arquitetura}
\begin{itemize}
  \item $n_x = 2$ entradas ($x_1, x_2$)
  \item $n_h = 2$ neurônios ocultos (ativação $\tanh$)
  \item $n_y = 1$ saída (ativação sigmoide)
  \item $\mathcal{L}$ = BCE
\end{itemize}

\subsubsection*{Forward Pass}
\begin{align}
  \mathbf{z}^{[1]} &= \mathbf{W}^{[1]}\mathbf{x} + \mathbf{b}^{[1]} \label{eq:fwd1} \\
  \mathbf{a}^{[1]} &= \tanh(\mathbf{z}^{[1]}) \label{eq:fwd2} \\
  z^{[2]} &= \mathbf{w}^{[2]T}\mathbf{a}^{[1]} + b^{[2]} \label{eq:fwd3} \\
  \hat{y} &= \sigma(z^{[2]}) \label{eq:fwd4}
\end{align}

\subsubsection*{Backward Pass (da saída para a entrada)}
\begin{align}
  \frac{\partial \mathcal{L}}{\partial z^{[2]}} &= \hat{y} - y \label{eq:bwd1} \\
  \frac{\partial \mathcal{L}}{\partial \mathbf{w}^{[2]}} &=
    \frac{\partial \mathcal{L}}{\partial z^{[2]}} \cdot \mathbf{a}^{[1]T} \label{eq:bwd2} \\
  \frac{\partial \mathcal{L}}{\partial \mathbf{a}^{[1]}} &=
    \mathbf{w}^{[2]} \cdot \frac{\partial \mathcal{L}}{\partial z^{[2]}} \label{eq:bwd3} \\
  \frac{\partial \mathcal{L}}{\partial \mathbf{z}^{[1]}} &=
    \frac{\partial \mathcal{L}}{\partial \mathbf{a}^{[1]}} \odot \tanh'(\mathbf{z}^{[1]}) \label{eq:bwd4} \\
  \frac{\partial \mathcal{L}}{\partial \mathbf{W}^{[1]}} &=
    \frac{\partial \mathcal{L}}{\partial \mathbf{z}^{[1]}} \cdot \mathbf{x}^T \label{eq:bwd5}
\end{align}

Para dados de treinamento $\{(0,0)\to0, (0,1)\to1, (1,0)\to1, (1,1)\to0\}$,
a rede aprende representações intermediárias no espaço oculto que tornam os
padrões linearmente separáveis.

\subsection{Backpropagation em CNNs}

Nas redes convolucionais (Capítulo 7), a retropropagação através das camadas
de convolução requer a operação de \textbf{convolução transposta} (também
chamada \textit{deconvolution} ou \textit{fractionally strided convolution}):

\begin{equation}
  \frac{\partial \mathcal{L}}{\partial \mathbf{X}} =
    \frac{\partial \mathcal{L}}{\partial \mathbf{Z}} \star_{\text{rot}} \mathbf{K}
  \label{eq:conv-backward}
\end{equation}

Onde $\star_{\text{rot}}$ é a correlação cruzada com o kernel rotacionado em
$180^\circ$ e $\mathbf{Z} = \mathbf{X} \ast \mathbf{K}$.

\destaque{
\textbf{Por que a backpropagation funciona?} O Teorema da Aproximação Universal
(Cybenko, 1989; Hornik et al., 1989) garante que uma MLP com \textbf{uma única
camada oculta} e ativação não-linear pode aproximar qualquer função contínua
em um compacto com precisão arbitrária --- desde que haja neurônios suficientes.
Na prática, redes profundas (múltiplas camadas) são exponencialmente mais
eficientes em número de parâmetros que redes rasas equivalentes.
}

%-----------------------------------------------------------------------------
\section{Otimizadores}
\label{sec:otimizadores}

O otimizador implementa a estratégia de atualização dos pesos a partir dos
gradientes calculados pela backpropagation. A escolha do otimizador impacta
diretamente a velocidade de convergência e a qualidade do mínimo encontrado.

\subsection{SGD (\textit{Stochastic Gradient Descent})}

O SGD atualiza os parâmetros usando o gradiente estimado em um \textbf{mini-batch}
de $m$ exemplos, ao invés do dataset completo (Batch GD) ou de um único exemplo
(SGD puro):

\begin{equation}
  \mathbf{W} \leftarrow \mathbf{W} - \eta \cdot \frac{1}{m} \sum_{i \in \text{batch}}
    \nabla_{\mathbf{W}} \mathcal{L}(\hat{y}^{(i)}, y^{(i)})
  \label{eq:sgd}
\end{equation}

\subsubsection{Hiperparâmetros do SGD}

\textbf{Hiperparâmetros do SGD:}
\begin{itemize}
  \item $\eta$ (learning rate): tipicamente $10^{-3}$ a $10^{-1}$
  \item \textbf{Batch size:} $32/64$ (CNNs), $128/256$ (MLPs com dados tabulares)
  \item \textbf{Momentum:} adiciona inércia $\mathbf{v} \leftarrow \beta\mathbf{v} + \eta\nabla\mathcal{L}$
        para amortecer oscilações ($\beta \approx 0.9$)
  \item \textbf{Nesterov:} calcula o gradiente no ponto ``antecipado''
        $\mathbf{W} - \beta\mathbf{v}$, melhorando a convergência em funções
        com curvatura acentuada
\end{itemize}

\subsection{Adam (\textit{Adaptive Moment Estimation})}

O Adam (Kingma \& Ba, 2015) combina as vantagens do RMSprop e do momentum,
sendo o otimizador padrão na maioria das aplicações odontológicas modernas
(\citeonline{DL-ODT-004}, \citeonline{ODT-009}).

\begin{align}
  m_t &= \beta_1 m_{t-1} + (1 - \beta_1) g_t \label{eq:adam-m} \\
  v_t &= \beta_2 v_{t-1} + (1 - \beta_2) g_t^2 \label{eq:adam-v} \\
  \hat{m}_t &= \frac{m_t}{1 - \beta_1^t} \label{eq:adam-mhat} \\
  \hat{v}_t &= \frac{v_t}{1 - \beta_2^t} \label{eq:adam-vhat} \\
  \theta_{t+1} &= \theta_t - \frac{\eta}{\sqrt{\hat{v}_t} + \epsilon} \hat{m}_t \label{eq:adam-update}
\end{align}

Onde $g_t = \nabla_{\theta} \mathcal{L}(\theta_t)$ é o gradiente no passo $t$,
$m_t$ é a média móvel do gradiente (primeiro momento), $v_t$ é a média móvel do
gradiente ao quadrado (segundo momento), e $\beta_1=0.9$, $\beta_2=0.999$,
$\epsilon=10^{-8}$ são os valores padrão.

\subsection{RMSprop}

Desenvolvido por Geoffrey Hinton, o RMSprop normaliza o gradiente pela raiz
quadrada da média móvel dos gradientes quadrados, sendo particularmente eficaz
para problemas com gradientes de magnitudes muito diferentes (comum em dados
de imagem odontológica com desbalanceamento severo):

\begin{align}
  v_t &= \rho v_{t-1} + (1 - \rho) g_t^2 \label{eq:rmsprop-v} \\
  \theta_{t+1} &= \theta_t - \frac{\eta}{\sqrt{v_t + \epsilon}} g_t \label{eq:rmsprop-update}
\end{align}

Com $\rho=0.9$ (fator de decaimento).

\begin{table}[H]
  \centering
  \caption{Comparação de otimizadores para treinamento de CNNs odontológicas}
  \label{tab:otimizadores}
  \begin{tabularx}{\textwidth}{lcccX}
    \toprule
    \textbf{Otimizador} & \textbf{LR típico} & \textbf{Momentum} & \textbf{Adaptativo} & \textbf{Uso em Odontologia} \\
    \midrule
    SGD + Momentum & $10^{-2}$ & 0.9 & Não & Fine-tuning ResNet (transfer learning) \\
    Adam & $10^{-4}$ & $\beta_1$=0.9 & Sim & Padrão para CNNs odontológicas \\
    AdamW & $10^{-4}$ & $\beta_1$=0.9 & Sim & Com weight decay desacoplado \\
    RMSprop & $10^{-3}$ & $\rho$=0.9 & Sim & Segmentação U-Net (Dice loss instável) \\
    \bottomrule
  \end{tabularx}
\end{table}

\notaexplicativa{O AdamW (Loshchilov \& Hutter, 2019) corrige uma falha do Adam
original: o weight decay é aplicado diretamente aos pesos ($\theta \leftarrow \theta - \eta\lambda\theta$)
ao invés de ser incorporado no gradiente adaptativo, resultando em melhor
generalização. Recomendamos AdamW para todos os projetos deste livro,
exceto quando houver justificativa específica.}

%-----------------------------------------------------------------------------
\section{Regularização e Combate ao Overfitting}
\label{sec:regularizacao}

O \textit{overfitting} ocorre quando o modelo memoriza o ruído e as idiossincrasias
do conjunto de treinamento, perdendo capacidade de generalização para dados não
vistos. Em odontologia, onde datasets são frequentemente pequenos (centenas a
poucos milhares de imagens), o overfitting é o \textbf{inimigo número um}.

\subsection{Regularização L1 e L2}

Adicionam penalidades à magnitude dos pesos na função de perda:

\begin{align}
  \mathcal{L}_{\text{total}} &= \mathcal{L}_{\text{data}} + \lambda R(\mathbf{W}) \label{eq:regularized-loss} \\
  R_{L2}(\mathbf{W}) &= \frac{1}{2} \sum_{l} \|\mathbf{W}^{[l]}\|_F^2 \quad \text{(Weight Decay)} \label{eq:l2} \\
  R_{L1}(\mathbf{W}) &= \sum_{l} \|\mathbf{W}^{[l]}\|_1 = \sum_{l} \sum_{i,j} |w_{ij}^{(l)}| \label{eq:l1}
\end{align}

\textbf{L2} penaliza pesos grandes quadraticamente, produzindo soluções ``suaves''
com todos os pesos pequenos mas não-nulos. \textbf{L1} produz \textbf{esparsidade}
(muitos pesos exatamente zero), atuando como seleção de features implícita.

Em odontologia, $\lambda \in [10^{-5}, 10^{-3}]$ é típico para CNNs.
\citeonline{DL-ODT-001} utilizou regularização L2 com $\lambda=10^{-4}$ em
seu ensemble CNN+Random Forest para classificação de condições dentárias,
reportando melhoria de 3.2\% na acurácia de validação.

\subsection{Dropout}

\subsubsection{Mecanismo do Dropout}

Dropout (Srivastava et al., 2014) é a técnica de regularização mais eficaz para
redes profundas. Durante o treinamento, cada neurônio tem probabilidade $p$ de
ser temporariamente ``desligado'' (sua saída é zerada):

\begin{equation}
  \tilde{a}_i^{(l)} = a_i^{(l)} \cdot m_i, \quad m_i \sim \text{Bernoulli}(1-p)
  \label{eq:dropout}
\end{equation}

Na \textbf{inferência}, todos os neurônios são usados, mas suas saídas são
escalonadas por $(1-p)$ para compensar a ausência de dropout:

\begin{equation}
  a_i^{(l),\text{test}} = (1-p) \cdot a_i^{(l)}
  \label{eq:dropout-test}
\end{equation}

\subsubsection{Taxas de Dropout Típicas}

\textbf{Taxas de dropout típicas em CNNs odontológicas:}
\begin{itemize}
  \item $p=0.5$ para camadas totalmente conectadas (fully connected)
  \item $p=0.1-0.3$ para camadas convolucionais (dropout espacial 2D)
  \item $p=0$ na camada de saída (nunca aplicar dropout na saída)
\end{itemize}

\subsection{Early Stopping}

A técnica mais simples e muitas vezes mais eficaz: interromper o treinamento
quando a perda de validação para de melhorar por $k$ épocas consecutivas
(paciência, tipicamente $k=10-20$).

\begin{equation}
  \text{Parar se } \mathcal{L}_{\text{val}}^{(t)} > \min_{j < t} \mathcal{L}_{\text{val}}^{(j)}
  \text{ por } k \text{ épocas consecutivas}
  \label{eq:early-stopping}
\end{equation}

\subsection{Batch Normalization}

\subsubsection{Mecanismo da Batch Normalization}

Batch Normalization (BN), proposta por Ioffe \& Szegedy (2015), normaliza as
ativações de cada camada para média zero e variância unitária durante o treinamento,
permitindo taxas de aprendizado mais altas e atuando como regularizador:

\begin{align}
  \mu_{\mathcal{B}} &= \frac{1}{m} \sum_{i=1}^{m} x_i \label{eq:bn-mean} \\
  \sigma_{\mathcal{B}}^2 &= \frac{1}{m} \sum_{i=1}^{m} (x_i - \mu_{\mathcal{B}})^2 \label{eq:bn-var} \\
  \hat{x}_i &= \frac{x_i - \mu_{\mathcal{B}}}{\sqrt{\sigma_{\mathcal{B}}^2 + \epsilon}} \label{eq:bn-norm} \\
  y_i &= \gamma \hat{x}_i + \beta \label{eq:bn-scale}
\end{align}

Os parâmetros $\gamma$ (escala) e $\beta$ (deslocamento) são aprendidos,
permitindo que a rede ``desfaça'' a normalização se necessário.

\subsubsection{Benefícios para CNNs Odontológicas}

\textbf{Benefícios da BN em CNNs odontológicas:}
\begin{enumerate}
  \item Permite $\eta$ de $10^{-2}$ a $10^{-1}$ (vs $10^{-4}$ sem BN)
  \item Reduz sensibilidade à inicialização dos pesos
  \item Efeito regularizador (ruído no mini-batch atua como dropout implícito)
  \item Acelera convergência em $14\times$ no paper original
\end{enumerate}

\destaque{
\textbf{Recomendação prática para o leitor:} Para qualquer CNN odontológica, a
combinação \textbf{BatchNorm + ReLU + AdamW + Early Stopping} é o ponto de
partida mais robusto. Acrescente Dropout ($p=0.5$ nas FC layers) se houver
sinais de overfitting. Regularização L2 ($\lambda=10^{-4}$) raramente prejudica
e frequentemente ajuda. Evite otimização prematura --- comece simples e adicione
complexidade conforme necessário.
}

%-----------------------------------------------------------------------------
\section{Arquitetura MLP para Diagnóstico Odontológico}
\label{sec:mlp-odonto}

Uma MLP típica para classificação de condições dentárias a partir de features
extraídas (idade do paciente, profundidade de sondagem, índice de placa,
mobilidade dentária, sangramento à sondagem) segue esta estrutura:

\begin{figure}[H]
  \centering
  \begin{tikzpicture}[
    node distance=1.5cm,
    neuron/.style={circle, draw, minimum size=1cm, fill=corSecundaria!20},
    input/.style={circle, draw, minimum size=1cm, fill=corPrimaria!20},
    output/.style={circle, draw, minimum size=1cm, fill=corDestaque!20},
    ]
    % Input layer
    \node[input] (I1) {$x_1$};
    \node[input, below=0.3cm of I1] (I2) {$x_2$};
    \node[below=0.3cm of I2] (DotsI) {$\vdots$};
    \node[input, below=0.3cm of DotsI] (In) {$x_n$};

    % Hidden layer 1
    \node[neuron, right=2cm of I1] (H11) {$h^{(1)}_1$};
    \node[neuron, right=2cm of I2] (H12) {$h^{(1)}_2$};
    \node[below=0.3cm of H12] (DotsH1) {$\vdots$};
    \node[neuron, right=2cm of In] (H1m) {$h^{(1)}_m$};

    % Hidden layer 2
    \node[neuron, right=2cm of H11] (H21) {$h^{(2)}_1$};
    \node[neuron, right=2cm of H12] (H22) {$h^{(2)}_2$};
    \node[below=0.3cm of H22] (DotsH2) {$\vdots$};
    \node[neuron, right=2cm of H1m] (H2k) {$h^{(2)}_k$};

    % Output layer
    \node[output, right=2.5cm of H21] (O1) {$y_1$};
    \node[output, right=2.5cm of H22] (O2) {$y_2$};
    \node[below=0.3cm of O2] (DotsO) {$\vdots$};
    \node[output, right=2.5cm of H2k] (Oc) {$y_C$};

    % Labels
    \node[above=0.5cm of I1, font=\footnotesize\bfseries] {Entrada};
    \node[above=0.5cm of H11, font=\footnotesize\bfseries] {Oculta 1};
    \node[above=0.5cm of H21, font=\footnotesize\bfseries] {Oculta 2};
    \node[above=0.5cm of O1, font=\footnotesize\bfseries] {Saída};

    % Bio
    \node[right=2cm of Oc, text width=3cm, font=\footnotesize] (Bio) {
      \textbf{Exemplo odontológico:}\\
      $x_1$: Prof. sondagem\\
      $x_2$: Índice de placa\\
      $x_3$: Mobilidade\\
      $\vdots$\\
      $y_i$: Diagnóstico
    };
  \end{tikzpicture}
  \caption{Arquitetura MLP genérica para diagnóstico odontológico com duas
  camadas ocultas. As features clínicas entram pela esquerda; as camadas ocultas
  aprendem representações intermediárias; a saída codifica o diagnóstico (binário
  ou multiclasse).}
  \label{fig:mlp-arch}
\end{figure}

\citacaoativa{doi-1111-odi-14789}{10.1111/odi.14789}{Redes neurais artificiais feedforward
demonstraram capacidade preditiva superior ao modelo TNM clássico para
predição de sobrevida em 5 anos de pacientes com carcinoma espinocelular
oral (AUC 0.89 vs 0.80) (Mafra et al., 2023).}

%-----------------------------------------------------------------------------
\section{Prática: MLP com \texttt{odontoia.models.cnn}}
\label{sec:pratica-mlp}

\niveltres

\begin{pratica}{Construção e Treinamento de MLP para Classificação de
Condições Dentárias}
  \textbf{Objetivo:} Implementar uma rede neural multicamadas (MLP) utilizando
  o módulo \texttt{odontoia.models.cnn.build\_simple\_cnn} para classificação
  de 4 condições dentárias (hígido, cárie, restauração, implante) a partir de
  features clínicas estruturadas.

  \textbf{Especificação:} SPEC-006.1 --- Classificador MLP Odontológico

  \textbf{Critérios de aceitação:}
  \begin{itemize}
    \item Acurácia de validação $\geq 82\%$
    \item F1-score macro $\geq 0.80$
    \item Tempo de treinamento $< 5$ minutos em GPU T4 (Colab)
    \item Curvas de perda (treino/validação) plotadas
    \item Matriz de confusão com 4 classes gerada
  \end{itemize}
\end{pratica}

\subsection{Preparação dos Dados}

Assumindo um dataset tabular com as seguintes features clínicas extraídas de
1.512 radiografias panorâmicas (\citeonline{DL-ODT-001}):

\begin{table}[H]
  \centering
  \caption{Features do dataset de classificação dentária}
  \label{tab:features-mlp}
  \begin{tabularx}{\textwidth}{llX}
    \toprule
    \textbf{Feature} & \textbf{Tipo} & \textbf{Descrição} \\
    \midrule
    idade\_paciente & Numérica (0-100) & Idade em anos \\
    densidade\_ossea & Categórica (D1-D4) & Classificação Lekholm \& Zarb \\
    posicao\_dente & Categórica (1-32) & Notação FDI \\
    profundidade\_sondagem & Numérica (0-15mm) & Profundidade de bolsa \\
    sangramento & Binária (0/1) & Sangramento à sondagem \\
    mobilidade & Ordinal (0-3) & Grau de mobilidade dentária \\
    indice\_placa & Numérica (0-100\%) & Índice de placa visível \\
    \bottomrule
  \end{tabularx}
\end{table}

\subsection{Implementação da MLP}

O módulo \texttt{odontoia.models.cnn.build\_simple\_cnn} fornece uma fábrica
de modelos com parâmetros configuráveis. Para construir uma MLP, especificamos
\texttt{model\_type='mlp'} (o padrão para dados tabulares, em contraste com
\texttt{'cnn'} para imagens):

\begin{lstlisting}[language=Python, caption={Construção e treinamento de MLP
com odontoia.models.cnn}, label={lst:mlp-train}]
import odontoia
from odontoia.models.cnn import build_simple_cnn
from odontoia.data.loader import load_dental_tabular
from sklearn.model_selection import train_test_split
from sklearn.preprocessing import StandardScaler, LabelEncoder
import numpy as np
import matplotlib.pyplot as plt
import seaborn as sns

# 1. Carregar dados tabulares
df = load_dental_tabular('datasets/condicoes_dentarias.csv')
X = df.drop('diagnostico', axis=1).values
y = df['diagnostico'].values

# 2. Pré-processamento
scaler = StandardScaler()
X = scaler.fit_transform(X)

le = LabelEncoder()
y_encoded = le.fit_transform(y)  # 0:higido, 1:carie, 2:restauracao, 3:implante

X_train, X_test, y_train, y_test = train_test_split(
    X, y_encoded, test_size=0.2, random_state=42, stratify=y_encoded
)

# 3. Construir modelo MLP
model = build_simple_cnn(
    input_shape=(X_train.shape[1],),     # n_features
    num_classes=4,                       # 4 condições dentárias
    model_type='mlp',                    # MLP, não CNN
    hidden_layers=[128, 64, 32],         # 3 camadas ocultas
    activation='relu',                   # ReLU nas ocultas
    output_activation='softmax',         # Softmax na saída
    dropout_rate=0.3,                    # Dropout 30%
    l2_reg=1e-4,                         # Regularização L2
    use_batch_norm=True                  # Batch Normalization
)

# 4. Compilar
model.compile(
    optimizer='adamw',
    learning_rate=1e-3,
    loss='sparse_categorical_crossentropy',
    metrics=['accuracy']
)

# 5. Treinar com early stopping
history = model.fit(
    X_train, y_train,
    validation_data=(X_test, y_test),
    epochs=200,
    batch_size=64,
    callbacks=[
        odontoia.callbacks.EarlyStopping(
            monitor='val_loss',
            patience=20,
            restore_best_weights=True
        ),
        odontoia.callbacks.ReduceLROnPlateau(
            monitor='val_loss',
            factor=0.5,
            patience=10
        )
    ],
    verbose=1
)

# 6. Avaliar
test_loss, test_acc = model.evaluate(X_test, y_test)
print(f"Acurácia no teste: {test_acc:.3f}")
\end{lstlisting}

\subsection{Visualização dos Resultados}

\begin{lstlisting}[language=Python, caption={Visualização de curvas de treino
e matriz de confusão}, label={lst:mlp-viz}]
# Curvas de perda e acurácia
fig, axes = plt.subplots(1, 2, figsize=(12, 4))

axes[0].plot(history.history['loss'], label='Treino', color='#1A5276')
axes[0].plot(history.history['val_loss'], label='Validação', color='#E67E22')
axes[0].set_title('Curva de Perda')
axes[0].set_xlabel('Época')
axes[0].set_ylabel('Perda (Cross-Entropy)')
axes[0].legend()
axes[0].grid(alpha=0.3)

axes[1].plot(history.history['accuracy'], label='Treino', color='#1A5276')
axes[1].plot(history.history['val_accuracy'], label='Validação', color='#E67E22')
axes[1].set_title('Curva de Acurácia')
axes[1].set_xlabel('Época')
axes[1].set_ylabel('Acurácia')
axes[1].legend()
axes[1].grid(alpha=0.3)

plt.tight_layout()
plt.savefig('mlp_curvas_treino.png', dpi=150, bbox_inches='tight')
plt.show()

# Matriz de confusão
from sklearn.metrics import confusion_matrix, classification_report

y_pred = np.argmax(model.predict(X_test), axis=1)
cm = confusion_matrix(y_test, y_pred)

plt.figure(figsize=(8, 6))
sns.heatmap(cm, annot=True, fmt='d', cmap='Blues',
            xticklabels=le.classes_,
            yticklabels=le.classes_)
plt.title('Matriz de Confusão — Classificador MLP')
plt.xlabel('Predito')
plt.ylabel('Real')
plt.tight_layout()
plt.savefig('mlp_matriz_confusao.png', dpi=150, bbox_inches='tight')
plt.show()

# Relatório de classificação
print(classification_report(y_test, y_pred, target_names=le.classes_))
\end{lstlisting}

\subsection{Resultados Esperados}

\begin{table}[H]
  \centering
  \caption{Métricas esperadas da MLP para classificação dentária (4 classes)}
  \label{tab:mlp-results}
  \begin{tabularx}{\textwidth}{lccccX}
    \toprule
    \textbf{Classe} & \textbf{Precisão} & \textbf{Recall} & \textbf{F1} & \textbf{Suporte} & \textbf{Nota} \\
    \midrule
    Hígido & 0.88 & 0.91 & 0.89 & 320 & Maior classe (baseline) \\
    Cárie & 0.84 & 0.81 & 0.82 & 280 & Confunde com hígido incipiente \\
    Restauração & 0.91 & 0.93 & 0.92 & 250 & Melhor classe (distinta) \\
    Implante & 0.86 & 0.83 & 0.84 & 210 & Confunde com restauração \\
    \midrule
    \textbf{Média Macro} & 0.87 & 0.87 & 0.87 & 1060 & \\
    \textbf{Acurácia Global} & \multicolumn{5}{c}{0.874} \\
    \bottomrule
  \end{tabularx}
\end{table}

\teste{TEST-006.1: Acurácia de validação $\geq 0.82$ -- APROVADO (0.874).
TEST-006.2: F1-score macro $\geq 0.80$ -- APROVADO (0.87).
TEST-006.3: Early stopping ativou corretamente na época 78, restaurando
melhores pesos da época 58. OVERFITTING CONTROL: divergência treino/val
$<$ 5\% após época 40.}

\notaexplicativa{O dataset de 1.512 radiografias com 11.137 anotações de
Golkarieh et al. (2025, \texttt{arXiv:2508.21088}) é um dos maiores datasets
públicos anotados para classificação de condições dentárias. O artigo original
reporta acurácia de 85.4\% com CNN+RF ensemble, consistente com nossos
resultados de MLP (87.4\%). A diferença deve-se ao uso de features
pré-extraídas de alta qualidade e Batch Normalization.}

%-----------------------------------------------------------------------------
\section{Síntese do Capítulo}
\label{sec:sintese-06}

Este capítulo estabeleceu os fundamentos matemáticos e práticos das redes neurais
artificiais, desde o perceptron de Rosenblatt até as MLPs modernas com BatchNorm,
Dropout e otimizadores adaptativos. Os conceitos aqui apresentados --- funções de
ativação, retropropagação, funções de perda, regularização --- são a base sobre a
qual as arquiteturas convolucionais do Capítulo 7 são construídas.

\textbf{Pontos-chave para o clínico-pesquisador:}
\begin{enumerate}
  \item \textbf{ReLU + BatchNorm + AdamW} é a combinação campeã para treinamento
        estável e rápido de redes neurais em dados odontológicos.
  \item \textbf{Early stopping} é a técnica de regularização mais negligenciada
        e mais eficaz --- implemente-a sempre.
  \item O \textbf{overfitting} é o principal risco em datasets odontológicos
        pequenos ($n < 1000$); Dropout (0.3--0.5) e L2 ($\lambda=10^{-4}$)
        são as primeiras linhas de defesa.
  \item A \textbf{backpropagation} não é caixa-preta: entender a regra da cadeia
        capacita o pesquisador a diagnosticar e corrigir problemas de treinamento
        (vanishing/exploding gradients).
  \item \textbf{Softmax + Cross-Entropy} é o par canônico para classificação
        multiclasse (ex.: 4 condições dentárias, 5 tipos de lesões bucais).
\end{enumerate}

No próximo capítulo, estenderemos estes conceitos para \textbf{redes neurais
convolucionais} (CNNs), a arquitetura que revolucionou o diagnóstico por imagem
em odontologia, permitindo que o modelo aprenda diretamente dos pixels --- sem
necessidade de extração manual de features.

\begin{flushright}
\itshape\small
--- ``All models are wrong, but some are useful.'' (George E. P. Box, 1976)
\end{flushright}

%==============================================================================
% FIM DO CAPÍTULO 6
%==============================================================================
